% Отчет

\makereporttitle
{Информатика и системы управления} % Название факультета
{Программное обеспечение ЭВМ и информационные технологии} % Название кафедры
{лабораторной работе \textnumero~1} % Название работы (в дат. падеже)
{Название курса} % Название курса (необязательный аргумент)
{Тема работы} % Тема работы
{} % Номер варианта (необязательный аргумент)
{Иванов~И.~И./ИУ7-00Б} % Номер группы/ФИО студента (если авторов несколько, их необходимо разделить запятой)
{Петров~П.~П.} % ФИО преподавателя

% % Расчетно-пояснительная записка к курсовой работе

% \makecourseworktitle
% {Информатика и системы управления} % Название факультета
% {Программное обеспечение ЭВМ и информационные технологии} % Название кафедры
% {Тема работы} % Тема работы
% {Иванов~И.~И./ИУ7-00Б} % Номер группы/ФИО студента (если авторов несколько, их необходимо разделить запятой)
% {Петров~П.~П.} % ФИО научного руководителя
% {Сидоров~С.~С.,Кузнецов~К.~К.} % ФИО консультанта (необязательный аргумент; если консультантов несколько, их необходимо разделить запятой)

% % Отчет по научно-исследовательской работе

% \makeresearchtitle
% {Информатика и системы управления} % Название факультета
% {Программное обеспечение ЭВМ и информационные технологии} % Название кафедры
% {Тема работы} % Тема работы
% {Иванов~И.~И./ИУ7-00Б} % Номер группы/ФИО студента (если авторов несколько, их необходимо разделить запятой)
% {Петров~П.~П.} % ФИО научного руководителя
% {Сидоров~С.~С.,Кузнецов~К.~К.} % ФИО консультанта (необязательный аргумент; если консультантов несколько, их необходимо разделить запятой)

% % Расчетно-пояснительная записка к выпускной квалификационной работе

% \makethesistitle
% {Информатика и системы управления} % Название факультета
% {Программное обеспечение ЭВМ и информационные технологии} % Название кафедры
% {Тема работы} % Тема работы
% {Иванов~И.~И./ИУ7-00Б} % Номер группы/ФИО студента (если авторов несколько, их необходимо разделить запятой)
% {Петров~П.~П.} % ФИО научного руководителя
% {Сидоров~С.~С.} % ФИО консультанта (необязательный аргумент; если консультантов несколько, их необходимо разделить запятой)
% {Кузнецов~К.~К.} % ФИО нормоконтролера
